\documentclass[12pt]{article}
\usepackage[T1]{fontenc}
\usepackage[utf8]{inputenc}
\usepackage{lmodern}
\usepackage{textcomp}
\usepackage{enumitem}
\usepackage{geometry}
\geometry{margin=1in}
\begin{document}
\begin{center}
Answer Key - Biology - Undergraduate Level Assessment 2 \
\small (Answers are ordered exactly as in the question paper.)
\end{center}

\section*{Multiple Choice}
\begin{enumerate}[label=\arabic*.]
\item \textbf{Answer:} C
\\ \textit{Options:} A) biomass of all producers; B) biomass of all consumers; C) energy available to heterotrophs; D) energy consumed by decomposers; E) total solar energy converted to chemical energy by producers
\\ \textit{Note:} Correct option: C - energy available to heterotrophs
\item \textbf{Answer:} A
\\ \textit{Options:} A) The stomach has a high permeability to hydrogen ions; B) The presence of beneficial bacteria that neutralize stomach acid; C) The stomach lining is replaced every month; D) The stomach lining secretes digestive enzymes only when food is present; E) The production of bile in the stomach neutralizes the digestive acids
\\ \textit{Note:} Correct option: A - The stomach has a high permeability to hydrogen ions
\item \textbf{Answer:} A
\\ \textit{Options:} A) mistakes in translation of structural genes.; B) mistakes in DNA replication.; C) translocations and mistakes in meiosis.; D) recombination at fertilization.
\\ \textit{Note:} Correct option: A - mistakes in translation of structural genes.
\item \textbf{Answer:} A
\\ \textit{Options:} A) that humans evolved from apes.; B) that evolution occurs.; C) that evolution occurs through a process of gradual change.; D) that the Earth is older than a few thousand years.; E) a mechanism for how evolution occurs.
\\ \textit{Note:} Correct option: A - that humans evolved from apes.
\item \textbf{Answer:} A
\\ \textit{Options:} A) Compare gene frequencies and genotype frequencies of two generations.; B) Evaluate the population's response to environmental stress; C) Assess the mutation rates of the population's genes; D) Compare the population size over generations; E) Calculate the ratio of males to females in the population
\\ \textit{Note:} Correct option: A - Compare gene frequencies and genotype frequencies of two generations.
\end{enumerate}

\section*{True / False}
\begin{enumerate}[label=\arabic*.]
\setcounter{enumi}{5}
\item \textbf{Answer:} True
\\ \textit{Options:} A) True; B) False
\\ \textit{Note:} LLM-generated true statement based on MCQ.
\item \textbf{Answer:} True
\\ \textit{Options:} A) True; B) False
\\ \textit{Note:} LLM-generated true statement based on MCQ.
\item \textbf{Answer:} True
\\ \textit{Options:} A) True; B) False
\\ \textit{Note:} LLM-generated true statement based on MCQ.
\item \textbf{Answer:} True
\\ \textit{Options:} A) True; B) False
\\ \textit{Note:} LLM-generated true statement based on MCQ.
\item \textbf{Answer:} False
\\ \textit{Options:} A) True; B) False
\\ \textit{Note:} LLM-generated false statement. Correct answer: 'Bundle sheath, phloem parenchyma, companion cell, sieve tube'.
\end{enumerate}

\section*{Long Form Response}
\begin{enumerate}[label=\arabic*.]
\setcounter{enumi}{10}
\item \textbf{Answer:} The first mammals originated during the Carboniferous Period of the Paleozoic Era.. Explain why this is the correct answer and provide supporting evidence. Reference key facts or concepts that support this choice.
\\ \textit{Note:} Long-form derived from MCQ; award credit for stating the correct idea and giving supporting reasoning.
\item \textbf{Answer:} Plant tissues can be fully characterized based on their evolutionary heritage. Explain why this is the correct answer and provide supporting evidence. Reference key facts or concepts that support this choice.
\\ \textit{Note:} Long-form derived from MCQ; award credit for stating the correct idea and giving supporting reasoning.
\end{enumerate}

\vfill
\noindent\textit{End of Answer Key}
\end{document}